\chapter{参数估计}
\begin{example}
设总体 \(X\sim N(\mu,4)\)，其中 \(\mu\) 未知，记 \(\theta=\mu\)，\(\Theta=\mathbb{R}\)。对总体作 \(4\) 次观测，得样本 \(X=(X_1,X_2,X_3,X_4)^{\mathrm T}\)。若观测值为 \(x=(8,9,11,10)^{\mathrm T}\)，则：
\begin{enumerate}
    \item \textbf{估计量：} 取样本均值作为 \(\mu\) 的点估计量，即 \(\hat{\mu}(X)=\bar X=\frac{1}{4}(X_1+X_2+X_3+X_4)\)。

    \item \textbf{估计值：} 将样本观测值代入估计量，得 \(\hat{\mu}(x)=\bar x=\frac{8+9+11+10}{4}=9.5\)。

    \item \textbf{置信区间：} 由于 \(\bar X\sim N\left(\mu,\frac{4}{4}\right)=N(\mu,1)\)，取置信水平 \(95\%\)，由 \(z_{0.975}=1.96\)，得 \(\mu\) 的 \(95\%\) 置信区间为 \([\bar X-1.96\frac{2}{\sqrt{4}},\ \bar X+1.96\frac{2}{\sqrt{4}}]=[\bar X-1.96,\ \bar X+1.96]\)。代入 \(\bar x=9.5\)，得 \([9.5-1.96,\ 9.5+1.96]=[7.54,11.46]\)。
\end{enumerate}
因此，本例中估计量为 \(\hat{\mu}(X)=\bar X\)，估计值为 \(\hat{\mu}(x)=9.5\)，\(95\%\) 置信区间为 \([7.54,11.46]\)。
\end{example}
\begin{definition}[样本中心矩]
设 \(X_1,X_2,\cdots,X_n\) i.i.d.，是抽自总体 \(X\) 的一个样本，则称 \(\hat{\alpha}_k\) 为总体 \(X\) 的 \(k\) 阶样本原点矩，\(\hat{\beta}_k\) 为总体 \(X\) 的 \(k\) 阶样本中心矩。
\end{definition}
\begin{note}
    \begin{itemize}
    \item 二阶样本中心矩：\(\hat{\beta}_2=\frac{1}{n}\sum_{i=1}^n (X_i-\bar X)^2\)，用来描述样本的离散程度，也就是数据围绕样本均值的波动大小。它与样本方差本质上是同一个东西，只差一个修正系数。

    \item 三阶样本中心矩：\(\hat{\beta}_3=\frac{1}{n}\sum_{i=1}^n (X_i-\bar X)^3\)，用来描述样本分布的偏斜程度。若它大于 \(0\)，通常表示分布右偏；若小于 \(0\)，通常表示分布左偏；若接近 \(0\)，通常表示分布较对称。

    \item 四阶样本中心矩：\(\hat{\beta}_4=\frac{1}{n}\sum_{i=1}^n (X_i-\bar X)^4\)，用来描述样本分布的峰态，反映数据分布是更尖、更平，或者尾部是否更厚。
\end{itemize}
\end{note}
\begin{note}
\(X\sim F_\theta(x)\) 表示随机变量 \(X\) 服从一个由参数 \(\theta\) 决定的分布，也就是说 \(X\) 的分布函数为 \(F_\theta(x)\)。其中 \(\theta\) 取不同的值，就对应这一族分布中的不同具体分布。F（x）可以是正态分布可以是均匀分布这个只是抽象的表示。
\end{note}
\begin{example}
设总体均值为 \(\mu\)，现不直接估计 \(\mu\)，而是要估计参数
\(\theta=\mu^2\)。令 \(\alpha=\mu\)，则 \(\theta\) 可写为
\(\theta=\theta(\alpha)=\alpha^2\)。又因为 \(\mu\) 的一个自然估计量是样本均值
\(\alpha(X)=\bar X\)，故由代入估计法，将 \(\theta(\alpha)=\alpha^2\) 中的 \(\alpha\) 用 \(\alpha(X)=\bar X\) 代替，得到 \(\theta\) 的估计量为
\[
\theta(\alpha(X))=(\bar X)^2.
\]
因此，\(\mu^2\) 的代入估计量为 \(\bar X^2\)。
\end{example}
\begin{dxtips}
    矩阵就是用来同时解方程的
\end{dxtips}
\begin{note}
上述方法称为\textbf{矩估计法}。其基本思想是：若总体分布中含有 \(k\) 个未知参数
\(\theta=(\theta_1,\theta_2,\dots,\theta_k)^{\mathrm T}\)，则选取总体的前 \(k\) 个原点矩
\(\alpha_1,\alpha_2,\dots,\alpha_k\)，并将它们表示为参数的函数，即
\(\alpha_l=\alpha_l(\theta_1,\theta_2,\dots,\theta_k)\), \(l=1,\dots,k\)。再用相应的样本原点矩
\(\frac{1}{n}\sum_{i=1}^n X_i,\frac{1}{n}\sum_{i=1}^n X_i^2,\dots,\frac{1}{n}\sum_{i=1}^n X_i^k\)
分别代替总体原点矩，得到关于 \(\theta_1,\theta_2,\dots,\theta_k\) 的方程组
\[
\begin{cases}
\alpha_1(\theta_1,\theta_2,\dots,\theta_k)=\frac{1}{n}\sum_{i=1}^n X_i,\\
\alpha_2(\theta_1,\theta_2,\dots,\theta_k)=\frac{1}{n}\sum_{i=1}^n X_i^2,\\
\qquad\vdots\\
\alpha_k(\theta_1,\theta_2,\dots,\theta_k)=\frac{1}{n}\sum_{i=1}^n X_i^k.
\end{cases}
\]
把该方程组的解 \(\hat\theta=(\hat\theta_1,\hat\theta_2,\dots,\hat\theta_k)^{\mathrm T}\) 作为参数 \(\theta\) 的估计量，称为\textbf{矩估计量}。
\end{note}

\begin{example}
设总体 \(X\sim N(\mu,\sigma^2)\)，其中 \(\mu,\sigma^2\) 都未知。现用矩估计法估计参数 \(\mu\) 与 \(\sigma^2\)。

总体的一阶原点矩与二阶原点矩分别为 \(E(X)=\mu\) 与 \(E(X^2)=\mu^2+\sigma^2\)。用样本原点矩 \(\bar X=\frac{1}{n}\sum_{i=1}^n X_i\) 和 \(\frac{1}{n}\sum_{i=1}^n X_i^2\) 分别代替它们，得到方程组
\[
\mu=\bar X,\qquad \mu^2+\sigma^2=\frac{1}{n}\sum_{i=1}^n X_i^2.
\]
解得
\[
\hat\mu=\bar X,\qquad \hat\sigma^2=\frac{1}{n}\sum_{i=1}^n X_i^2-\bar X^2.
\]
因此，\(\mu\) 与 \(\sigma^2\) 的矩估计量分别为 \(\bar X\) 和 \(\frac{1}{n}\sum_{i=1}^n X_i^2-\bar X^2\)。
\end{example}
\begin{definition}[似然函数likelihood function，要求互相独立服从一种类型的分布]
设 \(X_1,X_2,\cdots,X_n\) i.i.d.，\(X_1\sim f(x;\theta)\)，\(\theta\in\Theta\)。令
\[
L(\theta;x)=L(\theta;x_1,x_2,\cdots,x_n)=\prod_{i=1}^n f(x_i;\theta)
\]
为样本 \(X=(X_1,\cdots,X_n)^{\mathrm T}\) 的概率密度（或概率函数），视其为 \(\theta\) 的函数，称为相应于样本 \(X\) 的似然函数；也称 \(L(\theta;X)\) 为相应于样本 \(X\) 的似然函数。
\end{definition}
\begin{dxtips}
似然就是想通过这个函数确定参数的大小    
\end{dxtips}
\begin{definition}[MLE用来确定参数的]
设相应于样本 \(X=(X_1,X_2,\cdots,X_n)^{\mathrm T}\) 的似然函数为 \(L(\theta;x)\)。若 \(\hat{\theta}(x)=\hat{\theta}(x_1,x_2,\cdots,x_n)\) 满足
\[
L(\hat{\theta}(x);x)=\max_{\theta\in\Theta} L(\theta;x),
\]
则称 \(\hat{\theta}(X)\) 为 \(\theta\) 的最大似然估计量，简记为 MLE（Maximum likelihood estimator）。
\end{definition}
\begin{theorem}[求最大似然估计量的三种常见情形]
设似然函数为 \(L(\theta;x)\)，则求最大似然估计量时，常见有以下三种情形：
\begin{itemize}
    \item 若 \(L(\theta;x)\) 关于 \(\theta\) 连续可导，且 \(\Theta\) 为开区域，则通常先解似然方程组
    \[
    \frac{\partial L(\theta;x)}{\partial \theta_j}=0,
    \qquad j=1,2,\dots,k,
    \]
    或等价地解对数似然方程组
    \[
    \frac{\partial \ln L(\theta;x)}{\partial \theta_j}=0,
    \qquad j=1,2,\dots,k.
    \]

    \item 若 \(L(\theta;x)\) 关于 \(\theta\) 有间断点，则不能直接用似然方程，需结合似然函数的具体形式分段讨论，在定义域内直接比较其大小。

    \item 若参数 \(\theta\) 取离散值，则常通过比较相邻参数值处似然函数的大小，来确定使 \(L(\theta;x)\) 取得最大值的参数。
\end{itemize}
\end{theorem}
\begin{theorem}[最大似然估计在一一变换下会“跟着变过去”]
设 \(X\sim f(x;\theta)\), \(\theta\in\Theta\), \(X\) 是来自于该分布的一个样本，\(\hat{\theta}(X)\) 是 \(\theta\) 的最大似然估计量。设 \(\eta=g(\theta)\) 是由 \(\Theta\) 到 \(\hat{\Theta}\) 上的一一映射，则 \(g(\hat{\theta}(X))\) 是 \(\eta\) 的最大似然估计量。
\end{theorem}
\begin{note}
  可以让整个最大的θ可以不止一个
\end{note}
\begin{definition}[M估计]
设样本为 \(X_1,\dots,X_n\)，参数为 \(\theta\)。若 \(\hat{\theta}\) 是通过最小化
\[
\sum_{i=1}^n \rho(X_i,\theta)
\]
得到的估计，其中 \(\rho\) 是关于误差的某个对称凸函数，则称 \(\hat{\theta}\) 为参数 \(\theta\) 的一个 \textbf{M估计}。

等价地，若 \(\hat{\theta}\) 满足估计方程
\[
\sum_{i=1}^n \psi(X_i,\theta)=0,
\]
其中 \(\psi\) 常可看作某个函数 \(\rho\) 的导数，则由该方程得到的估计也称为 \textbf{M估计}。
\end{definition}
\begin{dxtips}
    很多 M 估计就是为了减弱离群点的影响
\end{dxtips}
\begin{example}
设样本为
\[
x_1=1,\quad x_2=2,\quad x_3=10,
\]
现用 M 估计来估计位置参数 \(\theta\)。

\begin{enumerate}
    \item 若取
    \[
    \rho(u)=u^2,
    \]
    则对应的 M 估计为最小化
    \[
    \sum_{i=1}^3 \rho(x_i-\theta)
    =(1-\theta)^2+(2-\theta)^2+(10-\theta)^2.
    \]
    解得
    \[
    \hat{\theta}=\frac{1+2+10}{3}=\frac{13}{3}.
    \]
    此时得到的就是样本均值。

    \item 若取
    \[
    \rho(u)=|u|,
    \]
    则对应的 M 估计为最小化
    \[
    \sum_{i=1}^3 \rho(x_i-\theta)
    =|1-\theta|+|2-\theta|+|10-\theta|.
    \]
    解得
    \[
    \hat{\theta}=2.
    \]
    此时得到的就是样本中位数。
\end{enumerate}

这说明：M 估计的核心思想是先选定损失函数 \(\rho\)，再令总损失最小；不同的 \(\rho\) 会导出不同的估计量。
\end{example}

\begin{definition}[Bayes 估计]
设样本 $X = (X_1, X_2, \ldots, X_n)^\mathrm{T}$ 来自条件分布 $f(x \mid \theta)$，参数 $\theta$ 的先验分布为 $\pi(\theta)$。在给定样本 $X$ 后，由贝叶斯公式，$\theta$ 的后验分布为
\[
  p(\theta \mid X) = \frac{f(X \mid \theta)\,\pi(\theta)}{\displaystyle\int_{\Theta} f(X \mid \theta)\,\pi(\theta)\,\mathrm{d}\theta}.
\]
对于待估参数函数 $g(\theta)$，若以其后验均值
\[
  \hat{g}(X) = \int_{\Theta} g(\theta)\,p(\theta \mid X)\,\mathrm{d}\theta
             = \frac{\displaystyle\int_{\Theta} g(\theta)\,f(X \mid \theta)\,\pi(\theta)\,\mathrm{d}\theta}
                    {\displaystyle\int_{\Theta} f(X \mid \theta)\,\pi(\theta)\,\mathrm{d}\theta}
\]
作为点估计，则称 $\hat{g}(X)$ 为 $g(\theta)$ 的 \textbf{Bayes 估计}，也称\textbf{后验均值估计}。
\end{definition}
\begin{dxtips}[$\propto$ 表示正比于]
    

（proportional to），$A \propto B$ 即 $A = c \cdot B$，其中 $c$ 为与关心变量无关的常数。在贝叶斯推断中，由于
\[
p(\mu|\boldsymbol{X}) = \frac{f(\boldsymbol{X}|\mu)\cdot\pi(\mu)}{f(\boldsymbol{X})}
\]
分母 $f(\boldsymbol{X})$ 与 $\mu$ 无关，故直接写
\[
p(\mu|\boldsymbol{X}) \propto f(\boldsymbol{X}|\mu)\cdot\pi(\mu)
\]
省去分母使计算更简洁，配方完成后自然归一化。
\end{dxtips}
\begin{note}
三个关键概念：
\begin{enumerate}
  \item \textbf{先验分布 $\pi(\theta)$}：在看到数据之前，你对 $\theta$ 的主观判断。比如你觉得一枚硬币正面概率大概在 $0.5$ 左右，这就是先验。
  \item \textbf{似然函数 $f(X\mid\theta)$}：给定参数 $\theta$，观测到样本 $X$ 的概率。这是数据给你的信息。
  \item \textbf{后验分布 $p(\theta\mid X)$}：看到数据 $X$ 之后，对 $\theta$ 的更新判断：
\[
  p(\theta \mid X) = \frac{f(X \mid \theta)\,\pi(\theta)}{\displaystyle\int_{\Theta} f(X \mid \theta)\,\pi(\theta)\,\mathrm{d}\theta}
\]
简单说就是：后验 $\propto$ 似然 $\times$ 先验，分母只是归一化常数让概率加起来等于 $1$。
\begin{dxtips}
    然后大公式里面的积分就是归一化常数和条件概率里面分母的累加是一样的这里是积分
\end{dxtips}
\end{enumerate}
\end{note}

\begin{example}[抛硬币的 Bayes 估计]
设正面概率 $\theta$ 的先验为 $\pi(\theta) \sim \text{Beta}(2,2)$，抛硬币 $10$ 次观测到 $7$ 次正面。由 Beta-Binomial 共轭，后验为
\[
  \theta \mid X \sim \text{Beta}(2+7,\ 2+3) = \text{Beta}(9,5),
\]
Bayes 估计（后验均值）为 $\hat{\theta} = \frac{9}{14} \approx 0.643$，相比频率派 MLE $\hat{\theta}_{\text{MLE}} = 0.7$，结果被先验适当拉向 $0.5$。
\end{example}
\section{3.3点估计量评价}
\begin{definition}{$g(\boldsymbol{\theta})$ 的无偏估计量}{def:unbiased}
设 $\hat{g}(\boldsymbol{X})$ 是 $g(\boldsymbol{\theta})$ 的估计量，若对于任意 $\boldsymbol{\theta} \in \Theta$ 都有
\[
E_{\boldsymbol{\theta}}(\hat{g}(\boldsymbol{X})) = g(\boldsymbol{\theta}),
\]
则称 $\hat{g}(\boldsymbol{X})$ 为 $g(\boldsymbol{\theta})$ 的\textbf{无偏估计量}。
\end{definition}
 估计量的期望等于真实值，就叫无偏估计
\begin{example}[样本均值是总体均值的无偏估计]
设总体均值 $\mu$ 未知，取样本 $X_1, X_2, X_3 = 4, 5, 9$，样本均值为
\[
  \bar{X} = \frac{4+5+9}{3} = 6.
\]
若真实总体均值 $\mu = 6$，则 $E(\bar{X}) = \mu = 6$，即 $\bar{X}$ 是 $\mu$ 的无偏估计。
\end{example}
\begin{note}
无偏性与有效性衡量的是\textbf{估计方法}的好坏，而非样本本身。对同一批样本，不同估计量的优劣由以下两个标准评价：
\begin{itemize}
  \item \textbf{无偏性}：估计量的期望等于真实参数，即无系统误差。
  \item \textbf{有效性}：在所有无偏估计量中，方差最小，即估计结果最稳定。
\end{itemize}
「样本取得好不好」对应的是\textbf{抽样方法}（如随机抽样、分层抽样），是独立的概念。
\end{note}
\begin{definition}[有效估计量]
设 $\hat{g}(\boldsymbol{X})$ 和 $\tilde{g}(\boldsymbol{X})$ 都是 $g(\theta)$ 的无偏估计量，若
\[
  \mathrm{Var}_\theta(\tilde{g}(\boldsymbol{X})) \geq \mathrm{Var}_\theta(\hat{g}(\boldsymbol{X})), \quad \forall \theta \in \Theta,
\]
且大于号至少在某个 $\theta \in \Theta$ 成立，则称 $\hat{g}(\boldsymbol{X})$ 比 $\tilde{g}(\boldsymbol{X})$ \textbf{有效}。比值
\[
  e = \frac{\mathrm{Var}_\theta(\hat{g}(\boldsymbol{X}))}{\mathrm{Var}_\theta(\tilde{g}(\boldsymbol{X}))}
\]
称为 $\tilde{g}(\boldsymbol{X})$ 相对于 $\hat{g}(\boldsymbol{X})$ 的\textbf{效率}。
\end{definition}
\begin{dxtips}
    g(X)就是对X操作算出一个θ然后，看见Var就知道他要比较稳定性了
\end{dxtips}
\begin{example}[样本均值与样本中位数的有效性比较]
设总体 $X \sim N(\mu, 1)$，取样本量 $n = 5$。样本均值 $\bar{X}$ 与样本中位数 $m$ 都是 $\mu$ 的无偏估计，但
\[
  \mathrm{Var}(\bar{X}) = \frac{1}{5} = 0.2, \qquad \mathrm{Var}(m) \approx \frac{\pi}{2 \times 5} \approx 0.314.
\]
效率为 $e = \dfrac{0.2}{0.314} \approx 0.637$，故 $\bar{X}$ 比 $m$ 更有效。
\end{example}
\begin{definition}{一致最小方差无偏估计量}{def:UMVUE}
设 $\hat{g}(\boldsymbol{X})$ 是 $g(\theta)$ 的一个无偏估计量，且 $\operatorname{Var}_{\theta}(\hat{g}(\boldsymbol{X})) < \infty$。
若对 $g(\theta)$ 的任意无偏估计量 $g^*(\boldsymbol{X})$，都有
\[
\operatorname{Var}_{\theta}(\hat{g}(\boldsymbol{X})) \leq \operatorname{Var}_{\theta}(g^*(\boldsymbol{X})), \quad \forall \theta \in \Theta,
\]
则称 $\hat{g}(\boldsymbol{X})$ 为 $g(\theta)$ 的\textbf{一致最小方差无偏估计量}，简记作 UMVUE
(Uniformly minimum variance unbiased estimator)。
\end{definition}
在所有无偏估计里，方差最小的那个就叫 UMVUE。
\begin{theorem}{UMVUE 的充要条件}{thm:UMVUE}
设 $\mathcal{U}$ 非空，则 $\delta(\boldsymbol{X}) \in \mathcal{U}$ 为 $g(\theta)$ 的 UMVUE 的充要条件是：
对任意 $\delta_0(\boldsymbol{X}) \in \mathcal{U}_0$，有
\[
E_{\theta}(\delta(\boldsymbol{X})\delta_0(\boldsymbol{X})) = 0, \quad \forall\, \theta \in \Theta.
\]
\end{theorem}
\begin{dxtips}
    U是所有无偏估计的集合U0就是期望是0的无偏估计的集合
\end{dxtips}
\begin{note}
符号说明：
\begin{itemize}
  \item $\mathcal{U}$：$g(\theta)$ 的所有无偏估计量构成的集合
  \item $\mathcal{U}_0$：所有零无偏估计量的集合，即满足 $E_\theta(\delta_0(\boldsymbol{X})) = 0$ 的估计量
  \item $\delta(\boldsymbol{X})$：待验证是否为 UMVUE 的候选估计量
\end{itemize}
定理的本质：$\delta(\boldsymbol{X})$ 是 UMVUE $\iff$ 它与所有零无偏估计量正交，即
\[
  E_\theta(\delta(\boldsymbol{X})\,\delta_0(\boldsymbol{X})) = 0, \quad \forall\,\theta \in \Theta,\ \forall\,\delta_0 \in \mathcal{U}_0.
\]
直觉上，$\delta_0$ 是「纯噪声」（期望为零，不含有效信息）。$\delta$ 与所有噪声正交，说明其中没有掺入任何多余成分，因此方差已达最小。
\end{note}
UMVUE 如果存在，则在以概率 1 相等的意义下是唯一的，这就是下面的定理。

\begin{theorem}{UMVUE 的唯一性}{thm:UMVUE_unique}
设 $\delta(\boldsymbol{X})$，$\delta'(\boldsymbol{X})$ 都是 $g(\boldsymbol{\theta})$ 的 UMVUE，则
\[
P_{\theta}(\delta(\boldsymbol{X}) = \delta'(\boldsymbol{X})) = 1, \quad \forall\, \theta \in \Theta.
\]
\end{theorem}
\begin{definition}{正则条件}{def:regular}
设分布族为 $\{f(x;\theta): \theta \in \Theta\}$，其中 $\theta$ 为一维未知参数，所谓正则条件是指：
\begin{enumerate}
    \item $\Theta$ 为开区间；
    \item 分布族有共同支撑，即支撑集 $A = \{x: f(x;\theta) > 0\}$ 与 $\theta$ 无关；
    \item $\forall x \in A$，$\frac{\partial f(x,\theta)}{\partial \theta}$ 在 $\Theta$ 上存在，且满足
    $E_{\theta}\!\left(\frac{\partial \ln f(x;\theta)}{\partial \theta}\right) = 0$，$\forall \theta \in \Theta$；
    \item Fischer 信息量
    \[
    I(\theta) = E_{\theta}\left[\left(\frac{\partial \ln f(x;\theta)}{\partial \theta}\right)^2\right]
    = \int_A \left(\frac{\partial \ln f(x;\theta)}{\partial \theta}\right)^2 f(x;\theta)\mathrm{d}x > 0, \quad \forall\, \theta \in \Theta.
    \]
\end{enumerate}
\end{definition}
\begin{note}
正则条件的本质是保证 C-R 下界的推导在数学上成立，每个条件各自排除一类麻烦：
\begin{itemize}
  \item \textbf{条件1（$\Theta$ 为开区间）}：保证对 $\theta$ 求导合法，排除参数在边界处导数无意义的情况。
  \item \textbf{条件2（支撑集与 $\theta$ 无关）}：排除积分上下限随 $\theta$ 变动的情况（如 $U(0,\theta)$），从而保证积分与求导可交换顺序。
  \item \textbf{条件3（得分函数期望为零）}：即 $E_\theta\!\left(\frac{\partial \ln f(x;\theta)}{\partial \theta}\right) = 0$，是积分与求导交换顺序的直接体现，也是推导 C-R 下界的关键一步。
  \item \textbf{条件4（$I(\theta) > 0$）}：保证 Fisher 信息量有限且为正，使得 C-R 下界 $\frac{1}{I(\theta)}$ 有实际意义。
\end{itemize}
四个条件合在一起，确保对数似然函数关于 $\theta$ 足够光滑，C-R 不等式方可成立。
\end{note}
\begin{note}[Fisher 信息量]
Fisher 信息量 $I(\theta)$ 度量样本 $X$ 中包含的关于参数 $\theta$ 的信息量，定义为
\[
I(\theta) = \operatorname{Var}\left(\frac{\partial}{\partial\theta}\ln f(X;\theta)\right).
\]

\textbf{为什么取对数？}
取对数是为了归一化：
\[
\frac{\partial}{\partial\theta}\ln f = \frac{1}{f}\cdot\frac{\partial f}{\partial\theta},
\]
相当于把导数除以 $f$ 本身，变成相对变化率，消除量纲影响。

\textbf{为什么取方差？}
$S(X;\theta) = \dfrac{\partial}{\partial\theta}\ln f(X;\theta)$ 仍是随机变量，方差是描述其平均波动程度最自然的选择。
又由 $E[S]=0$，故
\[
I(\theta) = \operatorname{Var}(S) = E(S^2),
\]
方差等于二阶矩，计算更方便。$I(\theta)$ 越大，$\theta$ 越好估计。
\end{note}
\begin{theorem}{Cramér-Rao 不等式}
设总体分布族 $\{f(x;\theta): \theta \in \Theta\}$ 满足正则条件 1—4，待估函数 $g(\theta)$ 有一阶导数
$g'(\theta)$。$\boldsymbol{X} = (X_1, X_2, \cdots, X_n)^{\mathrm{T}}$ 是抽自总体的一个样本，
$\hat{g}(\boldsymbol{X}) = \hat{g}(X_1, X_2, \cdots, X_n)$ 是 $g(\theta)$ 的无偏估计量，且满足
\[
\frac{\partial}{\partial \theta} \int_{A^n} \hat{g}(x_1, \cdots, x_n) \prod_{i=1}^n f(x_i; \theta)
\mathrm{d}x_1 \cdots \mathrm{d}x_n
= \int_{A^n} \hat{g}(x_1, \cdots, x_n) \frac{\partial}{\partial \theta}
\left(\prod_{i=1}^n f(x_i; \theta)\right) \mathrm{d}x_1 \cdots \mathrm{d}x_n,
\]
[需要满足积分顺序可交换]
则有
\begin{equation}
\operatorname{Var}_{\theta}(\hat{g}(\boldsymbol{X})) \geq \frac{(g'(\theta))^2}{nI(\theta)}, \quad \forall\, \theta \in \Theta.
\end{equation}
\end{theorem}
\begin{note}
该定理给出了任意无偏估计量方差的下界 $\frac{(g'(\theta))^2}{nI(\theta)}$，
称为 \textbf{C-R 下界}，无论构造多么巧妙的无偏估计量，其方差都不可能低于此值。下界由三部分决定：
\begin{enumerate}
    \item $(g'(\theta))^2$：待估函数本身的复杂程度，$g(\theta)$ 变化越剧烈，下界越大；
    \item $I(\theta)$：Fisher 信息量，反映总体模型本身所含的信息量，$I(\theta)$ 越大，
    参数越容易估计准确；$I(\theta)$ 可视作大小为 1 的样本所提供的信息量；
    \item $n$：样本量越大，下界越小，估计越精准；$nI(\theta)$ 称为样本的
    \textbf{Fisher 信息量}，反映样本 $(X_1, X_2, \cdots, X_n)^{\mathrm{T}}$ 所提供的总信息。
\end{enumerate}
若某无偏估计量的方差恰好等于 C-R 下界，则它已达到最优，必为 UMVUE，无法再改进；
若方差远大于下界，则说明仍有改进空间。
\end{note}

\section{估计量的大样本性质}
\begin{definition}[相合估计量与强相合估计量]
设 $\hat{g}_n = \hat{g}_n(X_1, X_2, \cdots, X_n)$ 为 $g(\theta)$ 的估计量。

若对任意固定的 $\theta \in \Theta$，对任给的 $\varepsilon > 0$，有
\[
\lim_{n \to \infty} P_\theta\{|\hat{g}_n - g(\theta)| > \varepsilon\} = 0,
\]
则称 $\hat{g}_n$ 为 $g(\theta)$ 的\textbf{相合估计量}，简记为 $\hat{g}_n \xrightarrow{P} g(\theta)$。

若对任何固定的 $\theta \in \Theta$ 都有
\[
P_\theta\left\{\lim_{n \to \infty} \hat{g}_n = g(\theta)\right\} = 1,
\]
则称 $\hat{g}_n$ 为 $g(\theta)$ 的\textbf{强相合估计量}，简记为 $\hat{g}_n \xrightarrow{a.s.} g(\theta)$。

强相合估计量必为相合估计量。
\end{definition}
\begin{itemize}
  \item \textbf{弱相合} $\Leftrightarrow$ 依概率收敛（$\xrightarrow{P}$）
  \item \textbf{强相合} $\Leftrightarrow$ 几乎处处收敛（$\xrightarrow{a.s.}$）almost surely 
\end{itemize}
\begin{theorem}[强相合估计量经过连续函数变换后仍然是强相合的。]
设 $G(\theta_1, \theta_2, \cdots, \theta_k)$ 在参数空间 $\Theta \subset \mathbf{R}^k$ 上连续，$\hat{T}_{in}$ 为 $\theta_i$ 的强相合估计量，$i = 1, 2, \cdots, k$，则 $G(\hat{T}_{1n}, \hat{T}_{2n}, \cdots, \hat{T}_{kn})$ 为 $G(\theta_1, \theta_2, \cdots, \theta_k)$ 的强相合估计量。
\end{theorem}
\begin{theorem}[矩估计量的强相合性]
设总体有直到 $k$（$k \geq 2$）阶的矩 $\alpha_i,\ i = 1,\cdots,k$，$\beta_j,\ j = 2,\cdots,k$。若 $g(\theta)$ 可表示为
\[
g(\theta) = G(\alpha_1, \alpha_2, \cdots, \alpha_k, \beta_2, \cdots, \beta_k)
\]
且 $G$ 为连续函数，记 $\hat{\alpha}_i$，$\hat{\beta}_j$ 分别为样本原点矩和样本中心矩，则
\[
G(\hat{\alpha}_1, \hat{\alpha}_2, \cdots, \hat{\alpha}_k, \hat{\beta}_2, \cdots, \hat{\beta}_k)
\]
为 $g(\theta)$ 的强相合估计量。
\end{theorem}
\begin{note}
矩估计量具有强相合性需满足三个条件：
总体有足够高阶的矩存在；$g(\theta)$ 能写成各阶矩的连续函数；用样本矩替换总体矩。满足这三条，矩估计量就是强相合的。
\end{note}
\begin{theorem}[最大似然估计量的存在性与相合性]
设分布族 $\{f(x;\theta) : \theta \in \Theta\}$ 满足：
\begin{enumerate}
    \item $\Theta$ 是有限集；
    \item 对于不同的参数值 $\theta$ 和 $\theta'$，所对应的分布不同；
    \item $\{f(x;\theta) : \theta \in \Theta\}$ 有共同支撑，即 $A = \{x : f(x;\theta) > 0\}$ 与 $\theta$ 无关。
\end{enumerate}
则对于简单随机样本 $(X_1, X_2, \cdots, X_n)^{\mathrm{T}}$，$\theta$ 的最大似然估计量 $\hat{\theta}_n = \hat{\theta}_n(X_1, X_2, \cdots, X_n)$ 存在，且 $\hat{\theta}_n$ 为 $\theta$ 的相合估计量。
\end{theorem}
\begin{note}
三个条件的直觉如下：
$\Theta$ 是有限集：$\theta$ 的取值范围不能无限大，保证最大值能取到；
不同参数对应不同分布：保证参数是"可识别的"，不然两个不同的 $\theta$ 长得一样，根本无法区分；
共同支撑与 $\theta$ 无关：保证似然函数对 $\theta$ 求导时不会出现支撑集随 $\theta$ 变化的麻烦，数学上处理起来干净。
\end{note}